% !PRINTED 707
\weberpage{707}
%<<CONT
Gattung ausdrücken, so kann die dritte Gattung nicht in diesem Ausdruck vorkommen, weil sonst die Residuen nicht $=0$ sein könnten. Dieser Ausdruck gibt also eine lineare Abhängigkeit zwischen den Differentialen zweiter Gattung, deren Unterecke $\mathfrak P^2, \mathfrak P_1^2, \mathfrak P_2^2, \ldots, \mathfrak P_p^2$ sind.

Nimmt man für irgend einen positiven Exponenten $m$:
\[
\mathfrak U = \mathfrak P^{m-1}\mathfrak P_1\mathfrak P_2\ldots\mathfrak P_p,
\]
und setzt
\[
z = \frac{\mathfrak A}{\mathfrak U}, \quad dz = \frac{\mathfrak Z}{\mathfrak U^2},
\]
so ist $\mathfrak Z$ durch $\mathfrak P^{m-2}$, $\mathfrak U^2$ durch $\mathfrak P^{2m-2}$, folglich das Untereck von $dz$ durch $\mathfrak P^m$ teilbar. Man bekommt dann eine lineare Abhängigkeit zwischen den Differentialen mit den Unterecken
\[
\mathfrak P^m, \mathfrak P^{m-1}, \ldots, \mathfrak P^2, \mathfrak P_1^2, \mathfrak P_2^2, \ldots, \mathfrak P_p^2,
\]
und damit ist der Satz 4. allgemein bewiesen.

\begin{center}
\rule{4cm}{0.4pt}
\end{center}
